\section{Quartic--sextic crossover}
\label{sec:crossover}

\subsection{Intrinsic crossover variable}

Define
\[
\boxed{
u=\frac{\delta}{\deltax}\ge0.
}
\]
Let
\[
R(u)=\frac{1+\sqrt{1+3u}}{2}.
\]
Then the scalar model gives
\begin{equation}
\boxed{
\frac{I_E}{I_0^{\rm scalar}}
=
R(u).
}
\label{eq:R-crossover}
\end{equation}

\subsection{Normalized action crossover}

Using the scalar residual action,
\begin{equation}
\boxed{
\frac{\mcJ_E}{\mcJ_0^{\rm scalar}}
=
4R^3-3R^2.
}
\label{eq:action-crossover}
\end{equation}

For \(u\gg1\), this reproduces the \(3/2\) action law. For \(u\to0\), the normalized scalar action tends to one.

\subsection{Effective action exponent}

Define
\[
\nu_{\mcJ}
=
\frac{d\log|\mcJ_E|}{d\log\delta}.
\]
The scalar model gives
\[
\boxed{
\nu_{\mcJ}
=
\frac{3(\sqrt{1+3u}-1)}
{2\sqrt{1+3u}-1}.
}
\]
Thus
\[
\nu_{\mcJ}\to\frac32
\quad(u\to\infty),
\qquad
\nu_{\mcJ}\to0
\quad(u\to0).
\]

\subsection{Leading Floquet crossover}

Near the diagonal,
\[
H_4
=
aI^2+4A\rho^2I^2x^2+\cdots.
\]
The physical second-return angle predicted by the scalar reduction is
\begin{equation}
\boxed{
\vartheta_F
=
8\rho I_E
\sqrt{
A\sqrt{a^2+3G\delta}
}.
}
\label{eq:thetaF}
\end{equation}

At \(\delta=0\),
\[
\vartheta_0^{\rm scalar}
=
8\rho I_0^{\rm scalar}\sqrt{A|a|}.
\]
Thus
\[
\boxed{
\frac{\vartheta_F}{\vartheta_0^{\rm scalar}}
=
R\sqrt{2R-1}.
}
\]

For \(u\gg1\),
\[
\vartheta_F\asymp\delta^{3/4}.
\]

\subsection{Parameter-free scalar relations}

Define
\[
X=\frac{I_E}{I_0^{\rm scalar}},
\qquad
Y=\frac{\mcJ_E}{\mcJ_0^{\rm scalar}},
\qquad
Z=\frac{\vartheta_F}{\vartheta_0^{\rm scalar}}.
\]
Within the scalar reduction,
\begin{equation}
\boxed{
Y=4X^3-3X^2,
}
\label{eq:YX}
\end{equation}
and
\begin{equation}
\boxed{
Z=X\sqrt{2X-1}.
}
\label{eq:ZX}
\end{equation}

These are coefficient-independent scalar crossover relations, not exact identities of the complete return-map germ.

\subsection{Corrections beyond the scalar reduction}

The complete sixth-order Hamiltonian is anisotropic:
\[
d\neq g,\qquad e\neq f.
\]
Therefore the true residual equilibrium is shifted slightly away from the exact diagonal. The scalar crossover curves should consequently be interpreted as leading normalized predictions. Their deviations measure angular sixth-order structure and higher-order terms.
